\documentclass[11pt]{article}
\usepackage[margin=1in]{geometry}
\usepackage{amsmath,amssymb}
\usepackage[hidelinks]{hyperref}

\title{Unbounded word lengths at fixed target arity\\[2pt]\large Disproof of Conjecture 00000002075}
\author{}
\date{}

\begin{document}
\maketitle

\noindent
For every fixed positive target degree $d$, morphisms from object $1$ to
object $d$ in the Lawvere theory of groups have unbounded shortest
representative lengths. This disproves the proposed $d^4$ bound and any
eventual polynomial bound in the target degree alone.

\section*{A family with exact word lengths}

Object $d$ is the $d$-fold power of the distinguished object. In the term
presentation of the group Lawvere theory, a map $m\to d$ is a $d$-tuple
of group words in $m$ variables, modulo the group identities. Equivalently,
it is a homomorphism $F_d\to F_m$, where $F_j$ is the free group on $j$
generators~\cite[p.~869]{lawvere}. Word length counts each generator letter
or inverse letter once, summed over the tuple. The empty
word represents the identity element $1$.

Fix $d\ge1$. For each nonnegative integer $N$ consider
\[
  f_N=(x^N,1,\ldots,1)\colon 1\longrightarrow d.
\]
Evaluate its terms in the additive group $\mathbb Z$, assigning the input
$x$ the value $1\in\mathbb Z$. The first coordinate of $f_N$ then has
value $N$. Any word representing that coordinate must have the same
value, since the group identities hold in $\mathbb Z$. Each letter
contributes $1$ or $-1$, so the length of the first word is at least $N$.
The representative with $N$ copies of $x$ in the first word and empty
words in all other coordinates attains that length. Thus the shortest
word length of $f_N$ is exactly $N$.

The argument applies at every fixed $d\ge1$. Given a proposed bound
$B(d)$, take an integer $N>B(d)$ at a positive target degree where the
bound is claimed.

\section*{Group terms and compositional expressions}

The same lower bound on variable occurrences holds if a representative
is written as a tree of terms in the usual group signature, with binary
multiplication, unary inverse, and the identity constant. Evaluated at
$x=1\in\mathbb Z$, a variable contributes $1$, an identity contributes
$0$, an inverse negates the value, and a multiplication adds its two
values. Induction and the triangle inequality show that the absolute
value of the evaluation is at most the number of variable occurrences.
Any term for the first coordinate still needs at least $N$ occurrences.
Even counting only multiplication and inverse symbols, a term tree with
$N$ variable occurrences has at least $N-1$ multiplication symbols,
because the number of leaves is one more than the number of binary nodes.

Explicit composition can shorten some expressions. For example, repeated
substitution of $x^2$ produces powers $x^{2^k}$ using $k$ copies of that
map. For any fixed finite alphabet used to record composition, symbols,
and indices, a length bound at a fixed $d$ permits only finitely many
strings. The maps $f_N$ remain pairwise distinct because their first
coordinates take different values at $1\in\mathbb Z$. This rules out a
target-only bound for that encoding as well. Providing a separate
unit-cost symbol for every power $x^N$ would instead introduce infinitely
many primitives.

\section*{Lean verification}

The Lean 4 project represents group terms and their interpretations in all
groups through the usual group operations and equations. It verifies
distinctness of the maps $f_N$, the lower bound on occurrences for any term
tuple representing $f_N$, and the exact word length. The theorem
\texttt{C2075.no\_eventual\_target\_bound} rules out any eventual target-only
bound on term occurrences. Run \texttt{lake build} in the \texttt{lean}
directory to check it.

\begin{thebibliography}{1}
\bibitem{lawvere}
F.~W. Lawvere, \emph{Functorial semantics of algebraic theories},
Proceedings of the National Academy of Sciences \textbf{50} (1963), 869--872.
\href{https://doi.org/10.1073/pnas.50.5.869}{DOI 10.1073/pnas.50.5.869}.
\end{thebibliography}

\end{document}
